\documentclass[../../../main.tex]{subfiles}

\begin{document}

\section{Algebra}

Continuing from the review on definitions and notations,
let's discuss some foundational knowledge in algebra that we learn
from competition math. I intend to keep this section short because
many of the topics reappear in olympiad style problems, and I think
we can discuss the generalization of the topics in latter notes.
In this section particular, let's discuss polynoimal expansion.

Before discussing binomial theorem, here is one notation to keep
in mind. I think many of us initially learn the notation for
combination as ${}_nC_k$. However in our notes, let's use the
standard notation $\binom{n}{k}$ as they are used for olympiad and
higher level math. Below is an important theorem to keep in mind for
our future proof.

\begin{theorem}{Pascal's Identity}{Pascal's Identity}
For positive integers $n > k$, the following equations hold.
\[
\binom{n}{k} = \binom{n-1}{k} + \binom{n-1}{k-1}
\]
\end{theorem}

There are few ways to prove this identity with easiest being drawing
Pascal's triangle in my opinion. However besides the visual proof,
here are two very famous proofs. In our first proof, we will be
approaching this problem with combinatoric ideas.

\begin{proof}
Assume that we are selecting $k$ leaders from $n$ people. You being
one of the people, there are two possibilities. In the first case,
you are the leader and the number of cases of choosing the rest of
the leaders is $\binom{n-1}{k-1}$. In the second case, you are not
a leader and there are $\binom{n-1}{k}$ number of cases of choosing
the remaining leaders. Therefore, $\binom{n}{k} = \binom{n-1}{k} +
\binom{n-1}{k-1}$.
\end{proof}

This was the first proof using combinatoric idea, and here is a proof
with a more algebraic approach.

\begin{proof}
Consider the following equations.
\begin{align*}
    \binom{n}{k}
    &= \frac{n!}{k! (n-k)!}
    = \frac{(n-1)! (n-k) + (n-1)! k}{k! (n-k)!} \\
    &= \frac{(n-1)!}{k! (n-k-1)!} + \frac{(n-1)!}{(k-1)! (n-k)!}
    = \binom{n-1}{k} + \binom{n-1}{k-1}
\end{align*}
Therefore, the identity holds.
\end{proof}

Personally, I prefer the first proof than the second one because it
feels more intuitive and second proof feels like a fallback just in
case I cannot think of combinatoric approach to an algebra problem,
but I think it depends on you!

Okay, with this in mind, let's discuss Binomial Theorem. I wanted to
discuss this because for some reason, it quite appears a lot in
unexpected problems and places like inequality and algebraic
combinatorics.

\begin{theorem}{Binomial Theorem}{Binomial Theorem}
For nonnegative integer $n$, the expansion of the term $(x + y)^n$
can be represented as the following.
\[
(x + y)^n = \sum_{k=0}^n \binom{n}{k} x^k y^{n-k}
\]
\end{theorem}

Again there are few different ways to prove this, and let's prove this
using two different ways.

\begin{proof}
Consider the term $x^k y^{n-k}$ for some appropriate integer $k$.
Notice that the number of such combination equals number of different
ways of choosing $k$ number of $(x + y)$ in $(x + y)^n$.
Therefore for all $x^k y^{n-k}$, their coefficient is $\binom{n}{k}$.
In other words, $(x + y)^n = \sum_{k=0}^n \binom{n}{k} x^k y^{n-k}$.
\end{proof}

Continuing with this combinatoric approach, here is another one using
induction.

\begin{proof}
First, notice that $x + y = \binom{1}{0} x + \binom{1}{1} y$.
Therefore, it suffices to show that the equation holds for $n + 1$
assuming that is holds for $n$. Consider the following equations.
\begin{align*}
    (x + y)^n (x + y)
    &= \left( \sum_{k=0}^n \binom{n}{k} x^k y^{n-k} \right) (x + y) \\
    &= \sum_{k=0}^n \left(
        \binom{n}{k} x^{k+1} y^{n-k} + \binom{n}{k} x^k y^{n-k+1}
    \right)
\end{align*}
Rewriting the terms, the following equation holds by Pascal's
Identity.
\begin{align*}
    (x + y)^n (x + y)
    &=\ \textcolor{Red}{\binom{n}{0} xy^n} + \binom{n}{0} y^{n+1} \\
    &\ \,+ \textcolor{Blue}{\binom{n}{1} x^2y^{n-1}} +
        \textcolor{Red}{\binom{n}{1} xy^n} \\
    &\ \,+ \binom{n}{2} x^3y^{n-2} +
        \textcolor{Blue}{\binom{n}{2} x^2y^{n-1}} \\
    &\qquad\vdots \\
    &=\ \binom{n}{n} x^{n+1} + \binom{n}{n} x^ny \\
    &= \binom{n}{0} y^{n+1} + \binom{n}{n} x^{n+1}
        + \sum_{k=1}^{n} \binom{n+1}{k} x^k y^{n+1-k} \\
    &= \sum_{k=0}^{n+1} \binom{n+1}{k} x^k y^{n+1-k}
\end{align*}
Therefore, the theorem holds.
\end{proof}

For this theorem too, I personally prefer the first method because
I think it is more intuitive and the second method of using induction
can works as a fallback. However, both proofs are both great,
and I think the choice depends on you!

This was a very short review of algebra that we learn from
competition math. Unlike geometry and number theory that requires
strong background understandings of topics from competition math,
I think algebra is one of those fields that are very different
when it comes to olympiad-style proof-based problems because there
are so many new topics that were not covered in competition math.
Although this section was short, we will discuss much in-depth in
future notes! For now, let's review some key properties from geometry.

\end{document}


